\documentclass[11pt]{article}

\usepackage[margin=1in]{geometry}
\usepackage{amsmath}
\usepackage{amssymb}
\usepackage{mathtools}
\usepackage{enumitem}
\usepackage{parskip}
\usepackage{hyperref}

\newcommand{\flops}{\text{FLOPs}}
\newcommand{\flopss}{\text{FLOPs/s}}
\newcommand{\bytesps}{\text{bytes/s}}
\newcommand{\AI}{\mathrm{AI}}
\newcommand{\Tmath}{T_{\text{math}}}
\newcommand{\Tmem}{T_{\text{mem}}}
\newcommand{\Tcomm}{T_{\text{comm}}}
\newcommand{\Tcomp}{T_{\text{comp}}}
\newcommand{\TPCIe}{T_{\text{PCIe}}}

\title{Scaling Book --- Section 2 Exercises}
\author{Rahul Bir}
\date{}

\begin{document}
\maketitle

\section*{Q1}

Sample from a 200B-parameter model in bf16. Split weights across 32 TPU v4p.

Per TPU:
\[
\frac{2 \cdot 200 \times 10^9 \text{ bytes}}{32} = 1.25 \times 10^{10} \text{ bytes/TPU}.
\]

TPU v4p HBM $\to$ VMEM bandwidth $= 1.2\mathrm{e}12$ bytes/s.

Total time to load to systolic array:
\[
\frac{1.25 \times 10^{10}}{1.2 \times 10^{12}} \,\text{s} \;\approx\; 0.01\,\text{s} \;=\; 10\,\text{ms}.
\]

If sampling from the model is memory bound, then this becomes a lower bound for the latency of sampling.

\section*{Q2}

\subsection*{(a) TPU v5e pod ($16 \times 16$ 2D torus)}

\begin{enumerate}[label=(\roman*)]
  \item Each host is connected to 2 four-chip trays:
  \[
  \text{total hosts} = \frac{16 \cdot 16}{8} = 32 \text{ cpu hosts}.
  \]
  \item Each v5e chip has 1 TPU core $\Rightarrow$ $256$ total TPU cores in the v5e pod.
  \item v5e FLOPs/s/chip (bf16) $= 1.97 \times 10^{14}$:
  \[
  \text{total pod } \flopss = 256 \cdot 1.97 \times 10^{14} \approx 5.04 \times 10^{16} \; \flopss.
  \]
  \item Total HBM $= 256 \times 16 \text{ GB} = 4096 \text{ GB}$.
\end{enumerate}

\subsection*{(b) TPU v5p pod ($16 \times 20 \times 28$ 3D torus)}

\begin{enumerate}[label=(\roman*)]
  \item One cpu host is connected to one 4-chip tray:
  \[
  \text{total hosts} = \frac{16 \times 20 \times 28}{4} = 2240 \text{ cpu hosts}.
  \]
  \item Each v5p chip has 2 TPU cores:
  \[
  \text{total cores} = 2 \cdot (16 \times 20 \times 28) = 17{,}920 \text{ TPU cores}.
  \]
  \item Total FLOPs/s (bf16):
  \[
  (16 \times 20 \times 28) \cdot 4.59 \times 10^{14} \approx 4.1 \times 10^{18} \; \flopss.
  \]
  \item Total HBM $= (16 \times 20 \times 28) \cdot 96 \text{ GB} = 860{,}160 \text{ GB} \approx 860 \text{ TB}$.
\end{enumerate}

\section*{Q3}

\[
A : \text{bf16}[D, F], \quad X : \text{bf16}[B, D], \quad X \cdot A : \text{bf16}[B, F].
\]
Assume $B \ll D$, $F = 4D$, on TPU v6e.

\begin{align*}
\text{total } \flops &= BF(D + D - 1) = 2BFD \\
\text{bytes moved} &= 2DF + 2BD + 2BF
\end{align*}

\textbf{Intensity (accel)} --- over PCIe:
\[
\text{Intensity}(\text{accel}) = \frac{9.20 \times 10^{14}}{1.6 \times 10^{10}} = 57{,}500.
\]

\textbf{Intensity (comp):}
\begin{align*}
\frac{2BFD}{2DF + 2BD + 2BF}
\;&\xrightarrow{F = 4D}\;
\frac{2B \cdot 4D \cdot D}{2 \cdot 4D \cdot D + 2BD + 2B \cdot 4D} \\
&= \frac{8BD^2}{8D^2 + 2BD + 8BD} \\
&= \frac{8BD^2}{8D^2 + 10BD}
\;\xrightarrow{B \ll D}\; B.
\end{align*}

$\therefore$ need $B \geq 57{,}500$ for the computation to remain FLOPs-bound over PCIe.

\section*{Q4}

\[
W : \text{int8}[16384, 4096], \quad x : \text{int8}[B, 4096], \quad x \cdot W \to \text{int8}[B, 16384] \quad \text{on 1 TPU v5e.}
\]

\textbf{Common quantities:}
\begin{align*}
\flops &= 16384 \cdot B \cdot (2 \cdot 4096) = 134{,}217{,}728\, B \approx 1.34 \times 10^8\, B \\
\text{bytes} &= (16384 \cdot 4096) + (B \cdot 4096) + (B \cdot 16384) \\
             &= 20{,}480\, B + 6.7 \times 10^7
\end{align*}

\subsection*{(i) HBM}
\begin{align*}
\Tmath &= \frac{\flops}{\text{v5e } \flopss} = \frac{1.34 \times 10^8\, B}{3.94 \times 10^{14}} \approx 3.4 \times 10^{-7}\, B \\
\Tmem &= \frac{\text{bytes}}{\text{v5e bw}} = \frac{20{,}480\, B + 6.7 \times 10^7}{8.2 \times 10^{11}}
\end{align*}
\[
\text{total time} \;\geq\; \max\{\Tmath,\, \Tmem\}.
\]

\textbf{Compute-bound threshold}\,\textsuperscript{$\dagger$}\textsuperscript{$\dagger$}\,\emph{Continuation from page 4 of the notes.}

\begin{align*}
\text{Intensity}(\text{comp}) &= \frac{1.34 \times 10^8\, B}{2 \times 10^4\, B + 6.7 \times 10^7} \\
\text{Intensity}(\text{accel}) &= \frac{3.94 \times 10^{14}}{8.2 \times 10^{11}} \approx 480
\end{align*}

To be compute-bound, $\text{Intensity}(\text{comp}) \geq \text{Intensity}(\text{accel})$:
\begin{align*}
(1.34 \times 10^8)\, B &\geq 480 \cdot (2 \times 10^4\, B + 6.7 \times 10^7) \\
(1.34 \times 10^8)\, B &\geq (9.6 \times 10^6)\, B + 480 \cdot 6.7 \times 10^7 \\
\left(1.34 \times 10^8 - 9.6 \times 10^6\right) B &\geq 480 \cdot 6.7 \times 10^7 \\
B &\geq \frac{480 \cdot 6.7 \times 10^7}{1.34 \times 10^8 - 9.6 \times 10^6} \approx 259.
\end{align*}

\subsection*{(ii) VMEM only}
VMEM bandwidth $\approx 22 \times$ HBM bandwidth, so
\[
\text{total time} \;\geq\; \max\!\left\{\, \underbrace{3.4 \times 10^{-7}\, B}_{\Tmath\text{ (stays same)}},\;\;
\underbrace{\frac{2 \times 10^4\, B + 6.7 \times 10^7}{22 \cdot 8.2 \times 10^{11}}}_{\Tmem\text{ (22$\times$ lower)}} \,\right\}.
\]

\textbf{Compute-bound threshold}\,\textsuperscript{$\dagger$}\textsuperscript{$\dagger$}\,\emph{Continuation from page 4 of the notes.}

Using VMEM only, $\text{Intensity}(\text{accel})$ drops by $22\times$, so
\[
B \geq \frac{259}{22} \approx 11.
\]

\section*{Q5}

TPU v5e $4 \times 4$ slice, per-hop latency is $1\,\mu\text{s}$. Want to send $\text{bf16}[8, 128, 8192]$ from $\text{TPU}\{0,0\}$ to $\text{TPU}\{3,3\}$.

Full v5e pod is $16 \times 16$. No wrap around connection in slice!

\emph{Note:} we have 2 output connections from $\text{TPU}\{0,0\}$ and 2 input connections to $\text{TPU}\{3,3\}$, so we can use both connections simultaneously to double our ICI bandwidth.

ICI b/w $= 9 \times 10^{10}$ (bidirectional through 2 paths).

We can think of each path like a pipe.

We have $6 \cdot 1\,\mu\text{s}$ startup time to fill the pipeline at which we reach steady state where data reaches by bytes/bw.

\begin{enumerate}[label=(\roman*)]
  \item Time for first byte to reach: $6 \cdot 1\,\mu\text{s} = 6\,\mu\text{s}$.
  \item Total time:
  \[
  6\,\mu\text{s} \;+\; \frac{2 \cdot 8 \cdot 128 \cdot 8192}{9 \times 10^{10}}\,\text{s}
  \;=\; \underbrace{6\,\mu\text{s}}_{\text{startup}} + \underbrace{186\,\mu\text{s}}_{\text{steady state}} = 192\,\mu\text{s}.
  \]
\end{enumerate}

\section*{Q6}

\[
A : \text{bf16}[128 \cdot 1024,\; 128 \cdot 1024], \quad \text{split across TPU v5e } 4 \times 4 \text{ slice but offloaded to host DRAM.}
\]

Want:
\begin{enumerate}[label=\arabic*.]
  \item Copy $A$ to $\text{TPU}\{0,0\}$.
  \item Multiply by vector $\text{bf16}[8,\; 128 \cdot 1024]$.
\end{enumerate}

\subsection*{Time for step 1 --- two options}

\textbf{Option 1.} PCIe-copy into TPU HBM, then transfer across ICI. (Note: ICI transfers happen from HBM, no need to move into VMEM first.)

\begin{itemize}[leftmargin=*]
  \item PCIe load time:
  \[
  \TPCIe \;=\; \frac{2 \cdot (128 \cdot 1024)^2}{16 \cdot 1.6 \times 10^{10}} \;=\; 0.07\,\text{s}.
  \]
  \item ICI transfer time. $\text{TPU}\{0,0\}$ has 2 input sockets, so effective ICI bw $= 2 \cdot 4.5 \times 10^{10}$ bytes/s. Neighbors are 1 hop away $\Rightarrow$ startup latency $1\,\mu\text{s}$. Afterwards, $\tfrac{15}{16}(128 \cdot 1024)^2$ bytes come in:
  \[
  \Tcomm \;=\; 1\,\mu\text{s} \;+\; \frac{\tfrac{15}{16}\,(128 \cdot 1024)^2}{2 \cdot 4.5 \times 10^{10}}\,\text{s} \;\approx\; 0.17\,\text{s}.
  \]
  \item Total time for step 1 (Option 1): $\Tcomm + \TPCIe \approx 0.24\,\text{s}$.
\end{itemize}

\textbf{Option 2.} From host DRAM, transfer using DCN.

\begin{itemize}[leftmargin=*]
  \item DCN bandwidth $= 6.25 \times 10^9$ bytes/s.
  \item Total time (also includes time to transfer from host DRAM to TPU HBM):
  \[
  \frac{\tfrac{15}{16} \cdot (128 \cdot 1024)^2}{6.25 \times 10^9} \,\text{s} \;=\; 2.57\,\text{s}.
  \]
\end{itemize}

\textbf{Option 1 is way better.}

\subsection*{Time for step 2 --- multiply by vector}
\begin{align*}
\Tmem  &= \frac{(128 \cdot 1024)^2 \;+\; 2 \cdot 8 \cdot 128 \cdot 1024 \;+\; 2 \cdot 8 \cdot 128 \cdot 1024}{8.2 \times 10^{11}} = 0.02\,\text{s} \\
\Tcomp &= \frac{2 \cdot (128 \cdot 1024)^2 \cdot 8}{1.97 \times 10^{14}} = 0.001\,\text{s}
\end{align*}

\subsection*{Bounds on total time}
\begin{align*}
\text{lower bound:} \quad \max\{\Tcomm,\; \Tmem,\; \Tcomp,\; \TPCIe\} &\approx 0.17 \,\text{s} \quad \text{(assuming perfect overlap)} \\
\text{upper bound:} \quad \Tcomm + \Tmem + \Tcomp + \TPCIe &\approx 0.26 \,\text{s}
\end{align*}

\appendix

\section*{Appendix A --- Pop Quiz}

\begin{itemize}[leftmargin=*]
  \item TPU v5p has 2 cores.
  \item Each core's VPU has $4 \cdot 8 \cdot 128$ ALUs.
\end{itemize}

\[
\text{TPU v5p vector } \flopss
\;=\;
\underbrace{(2 \cdot 4 \cdot 128 \cdot 8)}_{\flops}
\;\cdot\;
\underbrace{1.75\mathrm{e}9}_{1/\text{s}}
\]

\[
=\;
\left(2 \cdot 4 \cdot 128 \cdot 8 \cdot 1.75\mathrm{e}9\right)\,\flopss
\;=\;
1.43 \times 10^{13}\,\flopss.
\]

\end{document}
